\documentclass[11pt]{article}

\usepackage[margin=1in]{geometry}
\usepackage{amsmath,amssymb,mathtools,bm}
\usepackage{booktabs,longtable,array}
\usepackage[hidelinks]{hyperref}

\newcommand{\E}{\mathbb{E}}
\newcommand{\Cov}{\operatorname{Cov}}
\newcommand{\Var}{\operatorname{Var}}
\newcommand{\Law}{\operatorname{Law}}
\newcommand{\diag}{\operatorname{diag}}
\newcommand{\rank}{\operatorname{rank}}
\newcommand{\supp}{\operatorname{supp}}
\newcommand{\one}{\bm{1}}
\newcommand{\eps}{\boldsymbol{\epsilon}}
\newcommand{\etaV}{\boldsymbol{\eta}}
\newcommand{\ThetaV}{\boldsymbol{\Theta}}

\title{SCI-NOI-001 Independent Mathematical Core\\
Conditional sign ensembles and propagation}
\author{Independent Phase 1 audit derivation}
\date{Frozen against application parent
\texttt{d5015fe716971bf8ea617e8a187311bf5af05185}}

\begin{document}
\maketitle

\begin{abstract}
This document defines the conditional random ensemble generated by a
pre-mapmaking sign operation after the realized RTC/PTC state has been fixed.
It derives the ensemble mean, second moment, and covariance for arbitrary sign
correlation law; distinguishes randomization covariance from physical-noise
covariance; propagates the ensemble through fixed realized observation,
coaddition, validity, and filtering interfaces; and states the identity,
cardinality, dependence, and provenance required to interpret a finite stack.
It also defines a conceptual source-residual successor ensemble without
authorizing or assuming a fruit-loop implementation.

The scope is strictly SCI-NOI-001.  Sample-variance normalization, finite-count
correction, empirical variance or weight calibration, downstream weight and
signal-to-noise formulae, aperture-uncertainty estimators, thresholds,
feedback authority, and production realization-count/default policy belong
exclusively to SCI-NOI-002.  Thus ``distribution adequate'', ``covariance
adequate'', and ``variance-scale adequate'' below describe only the information
in an ensemble for a named use; they never validate a downstream estimator or
product.
\end{abstract}

\section{Authority, scope, and conditioning boundary}

This core uses only the approved abstract interfaces and frozen pre-core MAP
and FLT authority handoffs.  It makes no claim about the implementation at the
governing application SHA.  In particular, it does not inspect or infer the
actual randomized unit, sign sampler, seed derivation, balance rule,
realization count, mapmaker path, coaddition path, filter path, persistence
format, or fruit-loop state transition.  Those are Phase 2 trace questions.

For observation $o$, let $\mathcal I_o$ be the ordered set of sample--detector
coordinates admitted at the sign boundary.  An element of $\mathcal I_o$ has
an observation identity, scan and sample identity, detector UID, array/network
and map-group identity, and an upstream eligibility state.  Coordinates that
are missing, non-finite where finiteness is required, flagged out, or otherwise
ineligible are not silently represented by numerical zero in $\mathcal I_o$.
Zero, when present, is a valid numerical sample value.

Let $N_o=|\mathcal I_o|$.  The fixed realized post-RTC/PTC sample vector is
decomposed conceptually as
\begin{equation}
  x_o=s_o+n_o\in\mathbb R^{N_o},
  \qquad
  x_o\ \text{is fixed conditional on }\Theta_o .
  \label{eq:fixed-state}
\end{equation}
Here $s_o$ denotes deterministic astronomical or other deterministic content
and $n_o$ denotes the realized residual customarily called noise.  This
decomposition is for reasoning; neither term is asserted to be separately
observable.  The conditioning state $\Theta_o$ includes the realized RTC/PTC
outputs, residual correlations and imperfections, flags and eligibility,
detector coefficients, source-mask/selection state, observation membership,
geometry, map support, map operator, and every other fixed fact needed to make
the downstream operator deterministic.  SCI-NOI-001 conditions on those
facts; it does not validate their upstream correctness.

Let
\begin{equation}
  A_o:\mathbb R^{N_o}\longrightarrow\mathbb R^{P_o}
  \label{eq:map-operator}
\end{equation}
be the realized linear map operator admitted from the bounded MAP abstraction.
It may contain fixed gridding coefficients, normalization, geometry, and
support restrictions, but its admitted coefficient \texttt{weight\_I} is a
nonprecision gridding/normalization coefficient.  Nothing in this core
promotes it to variance, inverse variance, covariance, uncertainty, or
significance authority.  If an operational map step is data-dependent, then
Equation~\eqref{eq:map-operator} applies only after its exact realized operator
has been included in $\Theta_o$ and held fixed across the ensemble.

The state lifecycle is one-way:
\begin{equation}
  \text{requested}\longrightarrow\text{effective}
  \longrightarrow\text{observation-resolved}
  \longrightarrow\text{realized}.
  \label{eq:state-lifecycle}
\end{equation}
Realized RNG, assignment, operator, or product facts do not rewrite an earlier
request.

\section{Sign law and coherence units}

\subsection{Sample signs and a lifted coherence-unit law}

For realization $r$, let $\eps_o^{(r)}\in\{-1,+1\}^{N_o}$ and define
\begin{equation}
  D_{\eps_o^{(r)}}=\diag\!\left(\eps_o^{(r)}\right).
  \label{eq:sign-operator}
\end{equation}
The most useful general representation assigns one sign to each of $K_o$
disjoint coherence units.  Let
$c_o:\mathcal I_o\rightarrow\{1,\ldots,K_o\}$ be the unit map and let
$B_o\in\{0,1\}^{N_o\times K_o}$ have
$(B_o)_{ia}=1$ exactly when $c_o(i)=a$.  With unit signs
$\etaV_o^{(r)}\in\{-1,+1\}^{K_o}$,
\begin{equation}
  \eps_o^{(r)}=B_o\etaV_o^{(r)},
  \qquad
  (B_o\etaV_o^{(r)})_i=\eta_{o,c_o(i)}^{(r)}.
  \label{eq:coherence-lift}
\end{equation}
Thus a ``sign per sample'', ``sign per detector'', ``sign per scan'', or
``sign per observation'' is not mere metadata: it is a different partition
$c_o$ and therefore a different ensemble law.

For a random draw from the conditional sign law, define
\begin{align}
  m_{\eta,o} &= \E_{\eta}[\etaV_o\mid\Theta_o], &
  Q_{\eta,oo'} &= \E_{\eta}[\etaV_o\etaV_{o'}^{\mathsf T}
      \mid\ThetaV],
  \label{eq:unit-sign-moments}\\
  C_{\eta,oo'} &= Q_{\eta,oo'}-m_{\eta,o}m_{\eta,o'}^{\mathsf T},
  \label{eq:unit-sign-covariance}\\
  m_{\epsilon,o} &= B_o m_{\eta,o}, &
  Q_{\epsilon,oo'} &= B_oQ_{\eta,oo'}B_{o'}^{\mathsf T},
  \label{eq:sample-sign-moments}\\
  C_{\epsilon,oo'} &=Q_{\epsilon,oo'}
    -m_{\epsilon,o}m_{\epsilon,o'}^{\mathsf T}.
  \label{eq:sample-sign-covariance}
\end{align}
The matrix $Q_\epsilon$ is the sign-correlation law.  It is dimensionless,
has unit diagonal, is positive semidefinite for an attainable marginal law,
and must include cross-observation blocks when observations do not use
independent sign streams.  Centering, $m_\epsilon=0$, makes
$C_\epsilon=Q_\epsilon$, but a finite realized stack can have nonzero empirical
sign mean even under a centered target law.

For any fixed sample vector $v_o$, define its map contribution by coherence
unit,
\begin{equation}
  G_o(v_o)equiv
  \begin{bmatrix}g_{o1}(v_o)&\cdots&g_{oK_o}(v_o)\end{bmatrix},
  \qquad
  g_{oa}(v_o)=A_o\diag(v_o)B_oe_a .
  \label{eq:unit-map-contributions}
\end{equation}
Then $A_oD_{\epsilon_o}v_o=G_o(v_o)\etaV_o$.  This representation exposes
exactly which map modes a coherence choice preserves, destroys, or couples.

\subsection{Canonical sign-law limits}

The following are analytic laws, not assertions about the application.

\begin{enumerate}
\item Independent centered unit signs have
\begin{equation}
  m_{\eta,o}=0,\qquad Q_{\eta,oo}=I_{K_o}.
  \label{eq:independent-unit-law}
\end{equation}
Samples in one unit remain perfectly sign-coherent; samples in distinct units
have zero sign correlation.

\item One sign coherent over the whole observation has $K_o=1$ and
\begin{equation}
  Q_{\epsilon,oo}=\one\one^{\mathsf T}.
  \label{eq:global-coherent-law}
\end{equation}
It produces only $+A_ox_o$ and $-A_ox_o$ under a centered sign draw.

\item If $K_o$ is even and each realization is sampled uniformly from unit
assignments with exactly $K_o/2$ positive signs, then
\begin{equation}
  m_{\eta,o}=0,qquad
  (Q_{\eta,oo})_{ab}=
  \begin{cases}
    1,&a=b,\\
    -1/(K_o-1),&a\ne b,
  \end{cases}
  \quad
  Q_{\eta,oo}\one=0.
  \label{eq:within-realization-balance}
\end{equation}
Within-realization balance therefore changes the target covariance law; it is
not simply a way to obtain more accurate Monte Carlo centering.

\item A deterministic assignment $\etaV_o^\star$ has
\begin{equation}
  m_{\eta,o}=\etaV_o^\star,qquad
  Q_{\eta,oo}=\etaV_o^\star\etaV_o^{\star\mathsf T},
  \qquad C_{\eta,oo}=0.
  \label{eq:deterministic-sign-law}
\end{equation}
It has a nonzero raw second moment but no conditional randomization covariance.
\end{enumerate}

Balance across realizations is different from
Equation~\eqref{eq:within-realization-balance}.  Requiring
$\sum_r\eta_{oa}^{(r)}=0$ makes the empirical sign mean vanish coordinate by
coordinate but induces dependence among realizations.  Complement pairing,
orthogonal sign arrays, and sampling without replacement likewise alter the
joint law across $r$ even when every one-realization marginal has the desired
$Q_\eta$.  Both the within-realization law and the across-realization design
must therefore be stated.

\section{Current and conceptual residual ensembles}

The current conditional ensemble required by the dispatch is
\begin{equation}
  z_o^{(r)}
  =A_oD_{\epsilon_o^{(r)}}x_o
  =A_o\diag(x_o)\epsilon_o^{(r)}
  =G_o(x_o)\etaV_o^{(r)}.
  \label{eq:current-ensemble}
\end{equation}
The conceptual residual-state successor hypothesis is
\begin{equation}
  z_{o,\mathrm{res}}^{(r)}
  =A_oD_{\epsilon_o^{(r)}}[x_o-\widehat s_o],
  \label{eq:residual-ensemble}
\end{equation}
where $\widehat s_o$ is a fixed, realized, data-derived source model included
in the conditioning state.  Equation~\eqref{eq:residual-ensemble} is not an
authorized implementation, does not assert unbiasedness, and does not audit
the fruit loop.  It merely defines the mathematical object that would result
if the final source-subtracted sample state were randomized immediately before
a separately governed source re-addition.

\section{Conditional moments of the sign ensemble}

Write $X_o=\diag(x_o)$.  From
Equation~\eqref{eq:current-ensemble}, the conditional mean is
\begin{equation}
  \mu_o^{\epsilon}
  \equiv\E_\epsilon[z_o\mid\Theta_o]
  =A_oX_om_{\epsilon,o}
  =G_o(x_o)m_{\eta,o}.
  \label{eq:ensemble-mean}
\end{equation}
The uncentered ensemble second moment is
\begin{equation}
  M_{oo'}^{\epsilon}
  \equiv\E_\epsilon[z_oz_{o'}^{\mathsf T}\mid\ThetaV]
  =A_oX_oQ_{\epsilon,oo'}X_{o'}A_{o'}^{\mathsf T}
  =G_o(x_o)Q_{\eta,oo'}G_{o'}(x_{o'})^{\mathsf T}.
  \label{eq:ensemble-second-moment}
\end{equation}
The conditional randomization covariance is
\begin{equation}
  K_{oo'}^{\epsilon}
  \equiv\Cov_\epsilon(z_o,z_{o'}\mid\ThetaV)
  =A_oX_oC_{\epsilon,oo'}X_{o'}A_{o'}^{\mathsf T}
  =G_o(x_o)C_{\eta,oo'}G_{o'}(x_{o'})^{\mathsf T}.
  \label{eq:ensemble-covariance}
\end{equation}
Equivalently, before map projection,
\begin{equation}
  \Cov_\epsilon(D_\epsilon x\mid x,\Theta)
  =C_\epsilon\odot(xx^{\mathsf T}),
  \label{eq:hadamard-randomization}
\end{equation}
where $\odot$ is the Hadamard product.  Equations
\eqref{eq:ensemble-second-moment} and \eqref{eq:ensemble-covariance} are
distinct when the sign law is not centered.  A downstream calculation must
say which one it uses.

Replacing $x_o$ by $r_o=x_o-\widehat s_o$ in
Equations~\eqref{eq:ensemble-mean}--\eqref{eq:ensemble-covariance} gives the
moments of the conceptual residual ensemble.  No additional probability
claim follows from subtraction alone.

\section{Physical-noise covariance versus randomization covariance}

The only randomness in Equations~\eqref{eq:current-ensemble}--
\eqref{eq:ensemble-covariance} is the sign design: the realized $n_o$ is
fixed.  Expanding Equation~\eqref{eq:hadamard-randomization} gives
\begin{equation}
  C_{\epsilon,oo}\odot(x_ox_o^{\mathsf T})
  =C_{\epsilon,oo}\odot
  \left(s_os_o^{\mathsf T}+n_on_o^{\mathsf T}
  +s_on_o^{\mathsf T}+n_os_o^{\mathsf T}\right).
  \label{eq:fixed-realization-decomposition}
\end{equation}
Thus deterministic signal, the one realized noise outer product, and their
cross terms all enter the conditional randomization covariance.  Calling this
quantity ``the physical-noise covariance'' would be a category error without
additional conditions.

For comparison only, let $\Theta_o^-$ denote a hypothetical pre-realization
state and let
\begin{equation}
  \Sigma^{\mathrm{phys}}_{oo'}
  =\Cov(n_o,n_{o'}\mid\ThetaV^-)
  \label{eq:physical-noise-covariance}
\end{equation}
be the repeated-physical-noise covariance after the admitted upstream
processing.  The corresponding physical map covariance would be
\begin{equation}
  K^{\mathrm{phys,map}}_{oo'}
  =A_o\Sigma^{\mathrm{phys}}_{oo'}A_{o'}^{\mathsf T}.
  \label{eq:physical-map-covariance}
\end{equation}
Assume only for this comparison that physical noise has conditional mean zero
and is independent of a centered sign draw.  Averaging the sign-randomization
covariance over hypothetical physical repetitions yields
\begin{equation}
  \E_n[K_{oo'}^\epsilon\mid\ThetaV^-]
  =A_o\!left[
    C_{\epsilon,oo'}\odot
    \left(s_os_{o'}^{\mathsf T}+\Sigma^{\mathrm{phys}}_{oo'}\right)
  \right]A_{o'}^{\mathsf T}.
  \label{eq:average-randomization-covariance}
\end{equation}
The physical-noise part is therefore
\begin{equation}
  A_o[C_{\epsilon,oo'}\odot
  \Sigma^{\mathrm{phys}}_{oo'}]A_{o'}^{\mathsf T},
  \label{eq:preserved-physical-part}
\end{equation}
not Equation~\eqref{eq:physical-map-covariance} in general.  Independent unit
signs preserve within-unit physical covariance and suppress cross-unit
covariance.  A globally coherent sign preserves every covariance entry but
generates only the original map and its negative.  Exact within-realization
balance can introduce negative cross-unit sign correlations and annihilate a
common unit mode.  Covariance faithfulness is therefore a use- and
operator-specific equality or approximation statement, not a consequence of
having random signs.

A sufficient entrywise condition for preserving the physical covariance is
$C_{\epsilon,ij}=1$ wherever the relevant
$\Sigma^{\mathrm{phys}}_{ij}\ne0$, together with control of the deterministic
signal term.  It is not necessary: projection by $A$ or a named downstream
operator may make some discrepancies irrelevant.  The actual adequacy test
must be written on that projected domain and use an owner-approved criterion;
this core invents no tolerance.

\section{Deterministic-signal ensemble second-moment imprint}

Even for $m_\epsilon=0$, a fixed deterministic signal contributes
\begin{equation}
  M_{oo'}^{(s)}
  =A_o\diag(s_o)Q_{\epsilon,oo'}\diag(s_{o'})A_{o'}^{\mathsf T}
  =G_o(s_o)Q_{\eta,oo'}G_{o'}(s_{o'})^{\mathsf T}.
  \label{eq:signal-second-moment}
\end{equation}
This is the \emph{ensemble second-moment imprint}.  Zero ensemble mean does
not make it vanish.

The columns $g_{oa}(s_o)$ in
Equation~\eqref{eq:unit-map-contributions} give a morphology-based reading:
\begin{itemize}
\item For a compact source, many coherence-unit contributions can be aligned
around the same map location.  Independent unit signs give
$\sum_a g_{oa}g_{oa}^{\mathsf T}$, a positive localized second-moment imprint.
\item For a resolved or extended source, the same construction retains an
extended, potentially low-rank morphological imprint.  Random signs do not
turn deterministic structure into source-free noise.
\item For scan-synchronous deterministic structure, the imprint follows the
scan modes represented by $g_{oa}$.  Signs coherent on the same scan unit
preserve its internal mode, while signs independent across smaller units can
destroy cross-unit coherence.
\item For a globally coherent sign, the signal contribution is
$(A_os_o)(A_os_o)^{\mathsf T}$; the source is merely multiplied by $\pm1$.
For an exactly balanced unit law, any map mode represented by equal unit
columns is annihilated because $Q_\eta\one=0$.
\end{itemize}

A source-shaped elevation in a later empirical second-moment product can
therefore be a strong falsification diagnostic for deterministic-signal
leakage.  A downstream transformation of that elevation could manifest as a
source-shaped anti-weight symptom, but the weight definition and its validity
belong to SCI-NOI-002.  Absence of the symptom is not proof that
Equation~\eqref{eq:signal-second-moment} is zero: normalization, filtering,
support, finite ensemble sampling, or consumer behavior may hide it.

The complete fixed-realization second moment also contains the cross terms
\begin{equation}
  M^{(s,n)}_{oo'}=
  A_o\!left[
    Q_{\epsilon,oo'}\odot
    (s_on_{o'}^{\mathsf T}+n_os_{o'}^{\mathsf T})
  \right]A_{o'}^{\mathsf T},
  \label{eq:signal-noise-cross-imprint}
\end{equation}
which need not cancel in one realized data set.

\section{Conceptual fruit-residual bias trade}

To expose the trade without assuming a fruit-loop algorithm, write a purely
conceptual decomposition of the data-derived model as
\begin{equation}
  \widehat s_o=\widehat s_{o,s}+\widehat s_{o,n},
  \qquad
  r_o=x_o-\widehat s_o
  =(s_o-\widehat s_{o,s})+(n_o-\widehat s_{o,n}).
  \label{eq:residual-decomposition}
\end{equation}
The split into signal response $\widehat s_{o,s}$ and absorbed-noise response
$\widehat s_{o,n}$ is not generally observable; it is an analytic device.
With
$\delta s_o=s_o-\widehat s_{o,s}$ and
$n_{o,\mathrm{res}}=n_o-\widehat s_{o,n}$, the residual ensemble has signal
imprint
\begin{equation}
  M_{oo'}^{(\delta s)}
  =A_o\diag(\delta s_o)Q_{\epsilon,oo'}
   \diag(\delta s_{o'})A_{o'}^{\mathsf T}
  \label{eq:residual-signal-imprint}
\end{equation}
and a randomization-noise term formed from $n_{o,\mathrm{res}}$, not from
$n_o$.

Consequently, a good signal model can reduce the positive deterministic
imprint by making $\delta s_o$ small, while a data-derived model that absorbs
genuine noise can suppress physical-noise modes by making
$\widehat s_{o,n}$ nonzero.  The limiting cases are explicit:
\begin{align}
  \widehat s_o=s_o,\ \widehat s_{o,n}=0
  &\quad\Longrightarrow\quad r_o=n_o,
  \label{eq:perfect-signal-removal}\\
  \widehat s_o=x_o
  &\quad\Longrightarrow\quad r_o=0.
  \label{eq:complete-overfit}
\end{align}
Equation~\eqref{eq:perfect-signal-removal} is an oracle limit, not an
unbiasedness claim.  Equation~\eqref{eq:complete-overfit} shows why eliminating
source imprint alone cannot validate a residual ensemble.

If a later approved interface adds a fixed source map $b_o$ back to every
realization, then conditional covariance is unchanged while the uncentered
second moment and mean change:
\begin{equation}
  \Cov(b_o+z_{o,\mathrm{res}})=\Cov(z_{o,\mathrm{res}}),
  \qquad
  \E[(b_o+z)(b_o+z)^{\mathsf T}]
  =b_ob_o^{\mathsf T}+b_o\mu_z^{\mathsf T}
   +\mu_zb_o^{\mathsf T}+M_z.
  \label{eq:fixed-source-readdition}
\end{equation}
A realization-dependent or separately filtered re-addition does not satisfy
this fixed-offset identity and remains a future SCI-FRUIT-001 interface
question.

\section{Observation coaddition and cross-observation terms}

Let $H_o:\mathbb R^{P_o}\rightarrow\mathbb R^{P_c}$ be the exact realized,
fixed observation-to-coadd operator, including admitted membership, centered
integer embedding, fixed coefficient weighting, and fixed normalization.
Its coefficients remain nonprecision MAP coefficients.  The conditional coadd
ensemble is
\begin{equation}
  y^{(r)}=\sum_{o\in\mathcal O}H_oz_o^{(r)}.
  \label{eq:fixed-coadd}
\end{equation}
Therefore
\begin{align}
  \mu_y^\epsilon
  &=\sum_oH_o\mu_o^\epsilon,
  \label{eq:coadd-mean}\\
  M_y^\epsilon
  &=\sum_{o,o'}H_oM_{oo'}^\epsilon H_{o'}^{\mathsf T},
  \label{eq:coadd-second-moment}\\
  K_y^\epsilon
  &=\sum_{o,o'}H_oK_{oo'}^\epsilon H_{o'}^{\mathsf T}.
  \label{eq:coadd-covariance}
\end{align}
The $o\ne o'$ terms are controlled by the cross-observation sign law
$Q_{\epsilon,oo'}$ and by shared deterministic or realized-noise structure.
Independent centered sign streams make
$C_{\epsilon,oo'}=0$ for $o\ne o'$, eliminating conditional randomization
cross-covariance.  Reused seeds, shared assignments, complement designs, or a
global balance constraint can make those terms nonzero.  Observation-varying
geometry, valid region, membership, and coefficients are already represented
by the distinct $A_o$ and $H_o$; no universal scalar coadd guarantee follows.

If a coadd coefficient or membership is recomputed from the realizations,
write instead
\begin{equation}
  y^{(r)}=\sum_o H_o^{(r)}\!\left(z^{(1:R)},\ThetaV\right)z_o^{(r)}.
  \label{eq:data-derived-coadd}
\end{equation}
Equations~\eqref{eq:coadd-mean}--\eqref{eq:coadd-covariance} then do not apply:
the joint law contains weight--map and selection--map coupling.  SCI-NOI-001
requires that coupling to be identified.  It does not select, calibrate, or
validate the weight formula.  If coefficients are estimated once from the
original data and then held fixed across signs, they belong in $\ThetaV$ and
the fixed-operator equations apply conditionally; that conditioning still
does not establish physical precision.

\section{Filtering, support, and validity propagation}

For a fixed realized linear filtering operator $F$ admitted by the FLT
interface, define
\begin{equation}
  u^{(r)}=Fy^{(r)}.
  \label{eq:fixed-filter}
\end{equation}
Then
\begin{equation}
  \mu_u=F\mu_y,
  \qquad M_u=FM_yF^{\mathsf T},
  \qquad K_u=FK_yF^{\mathsf T}.
  \label{eq:filter-propagation}
\end{equation}
The same exact realized operator identity, support rule, response, and
provenance must be used for every realization and for any signal/kernel
quantity with which it is compared.  A realization-dependent $F^{(r)}$ must
instead be treated jointly with $y^{(r)}$; replacing it by one nominal matrix
would be an unsupported covariance claim.

Validity has its own domain and is not inferred from finite numerical output.
The contract requires:
\begin{enumerate}
\item the eligible input set $\mathcal I_o$ and its ordering remain fixed for
all realizations unless a different law explicitly includes eligibility as a
random variable;
\item multiplying by $\pm1$ does not clear an input flag, create support, or
turn a missing/non-finite value into an eligible value;
\item observation admission, map support, coadd membership, and filter-local
support are recorded separately;
\item a downstream finite value cannot promote an upstream-invalid sample or
pixel; and
\item every product retains raw-parent identity separately from local
numerical support and output validity.
\end{enumerate}
Direct pixelwise realization variance or signal-to-noise remains diagnostic
under the frozen FLT handoff.  Empirical-estimator and product validity are
reserved to SCI-NOI-002.

\section{Finite ensemble, dependence, and effective assignments}

Let the realized stack contain $R$ completed assignments.  Define its
empirical sign moments without choosing any downstream variance estimator:
\begin{align}
  \widehat m_{\epsilon,o}
  &=\frac{1}{R}\sum_{r=1}^{R}\epsilon_o^{(r)},
  \label{eq:empirical-sign-mean}\\
  \widehat Q_{\epsilon,oo'}
  &=\frac{1}{R}\sum_{r=1}^{R}
    \epsilon_o^{(r)}\epsilon_{o'}^{(r)\mathsf T}.
  \label{eq:empirical-sign-correlation}
\end{align}
The realized uncentered map second moment is exactly
\begin{equation}
  \frac{1}{R}\sum_{r=1}^{R}z_o^{(r)}z_{o'}^{(r)\mathsf T}
  =A_oX_o\widehat Q_{\epsilon,oo'}X_{o'}A_{o'}^{\mathsf T}.
  \label{eq:finite-stack-second-moment}
\end{equation}
The unnormalized centered scatter is
\begin{equation}
  \sum_{r=1}^{R}(z_o^{(r)}-\bar z_o)
  (z_{o'}^{(r)}-\bar z_{o'})^{\mathsf T}
  =A_oX_o\!left[
    \sum_{r=1}^{R}(\epsilon_o^{(r)}-\widehat m_{\epsilon,o})
    (\epsilon_{o'}^{(r)}-\widehat m_{\epsilon,o'})^{\mathsf T}
  \right]X_{o'}A_{o'}^{\mathsf T}.
  \label{eq:finite-stack-scatter}
\end{equation}
No divisor appears in Equation~\eqref{eq:finite-stack-scatter}: selecting a
sample-variance normalization or finite-$R$ correction is explicitly
SCI-NOI-002 work.

For $K$ unconstrained binary units, the assignment space has cardinality
\begin{equation}
  |\Omega_K|=2^K.
  \label{eq:assignment-cardinality}
\end{equation}
For exact within-realization half balance with even $K$,
\begin{equation}
  |\Omega_K^{\mathrm{bal}}|={K\choose K/2}.
  \label{eq:balanced-cardinality}
\end{equation}
Since $\etaV$ and $-\etaV$ yield $z$ and $-z$, they are distinct for a full
distribution but identical for every even statistic such as $zz^{\mathsf T}$.
For $K>0$, quotienting an unconstrained complement-closed assignment set by
this equivalence halves its second-moment pattern count.

For a realized stack define
\begin{align}
  U_{\mathrm{raw}}
  &=\#\{\etaV^{(r)}:1\le r\le R\},
  \label{eq:raw-unique-count}\\
  U_{\pm}
  &=\#\{\{\etaV^{(r)},-\etaV^{(r)}\}:1\le r\le R\}.
  \label{eq:complement-unique-count}
\end{align}
Duplicates reduce $U_{\mathrm{raw}}$; complement pairs reduce the number of
distinct even-moment patterns to at most $U_{\pm}$.  Neither count alone is a
universal statistical effective sample size because dependence is
use-specific.

When the joint randomization design across realizations is probabilistically
defined, let $h_r=h(z^{(r)})$ be a scalar summary with common marginal variance
$\sigma_h^2$ and cross-realization covariance matrix $\Gamma_h$.  A
use-specific effective size for the mean of that summary is
\begin{equation}
  R_{\mathrm{eff}}(h)
  \equiv
  \frac{R^2\sigma_h^2}{\one^{\mathsf T}\Gamma_h\one},
  \qquad
  \Var(\bar h)=\frac{\sigma_h^2}{R_{\mathrm{eff}}(h)},
  \label{eq:use-specific-effective-size}
\end{equation}
when the denominator is positive.  Complement pairs can be antithetic for an
odd $h$ but perfectly redundant for an even $h$; hence no single
$R_{\mathrm{eff}}$ describes all consumers.  A deterministic sign array has
no stochastic $R_{\mathrm{eff}}$ unless a randomization law over arrays is
also supplied.  Equation~\eqref{eq:use-specific-effective-size} is an
information/dependence descriptor only, not a finite-$R$ correction or a
production-count rule.

\section{RNG, identity, provenance, and reproducibility contract}

A seed value alone does not identify a sign ensemble.  Exact reproducibility
requires a scheduling-independent mapping
\begin{equation}
  (\text{algorithm/version},\ \text{master seed},\ \text{derivation rule},
  \ \text{observation ID},\ \text{iteration ID},\ \text{realization ID},
  \ \text{unit ID})
  \longmapsto\eta_{oa}^{(r)}.
  \label{eq:rng-key-map}
\end{equation}
The mapping may use independent substreams or a counter/key construction; the
mathematical requirement is that sequential/OpenMP scheduling, task arrival
order, and writer order cannot silently change an assignment.  Any intentional
shared stream or cross-observation coupling must be explicit in
$Q_{\eta,oo'}$ and provenance.

At minimum, the four lifecycle layers must carry the following facts:
\begin{longtable}{>{\raggedright\arraybackslash}p{0.16\textwidth}
                  >{\raggedright\arraybackslash}p{0.77\textwidth}}
\toprule
Layer & Required SCI-NOI-001 facts \\
\midrule
\endhead
Requested & Requested activation, realization count, sign/coherence policy,
balance or pairing request, seed mode/value if supplied, requested product
scope, and disabled expert values without rewriting them. \\
Effective & Effective activation and count; normalized sign-law identity;
RNG and seed policy; assignment sampler; within-realization and
across-realization constraints; complement/duplicate policy; and explicit
absence reasons for unsupported combinations. \\
Observation-resolved & Observation identity and order; eligible
sample--detector index and digest; detector UID and scan/sample indexing;
coherence partition and $K_o$; map/group/Stokes identity; geometry, WCS,
support and validity identity; operator identity; and observation-specific
assignment-space cardinality. \\
Realized & Application/config/input identities; RNG algorithm and version;
master seed and derivation; stream keys; exact realization and unit IDs;
assignment or assignment-digest inventory; completed count; failures and
missing realizations; $U_{\mathrm{raw}}$, $U_{\pm}$, balance, duplicates and
complements; empirical $\widehat m_\epsilon$ and $\widehat Q_\epsilon$ (or a
lossless digest-bound representation); exact observation/coadd/filter
operators; sequential/OpenMP policy; product inventory, shapes, units,
support, validity, and raw-parent provenance. \\
\bottomrule
\end{longtable}

The realization identity must be stable under persistence and joins.  A
minimal logical key is
\begin{equation}
  (\text{reduction},\text{iteration},\text{observation},
   \text{realization},\text{map/group},\text{component}),
  \label{eq:realization-identity}
\end{equation}
augmented by exact operator and assignment provenance.  A file position or
writer arrival order is not identity.  Requested, effective,
observation-resolved, and realized counts remain distinct; a partially
persisted stack cannot claim its requested cardinality.

Units and shape follow the operator chain.  $\epsilon$, $Q_\epsilon$, and
$C_\epsilon$ are dimensionless; $z_o$ has the map signal unit associated with
$A_o$; and $M$ and $K$ have squared signal unit.  Array, network, detector UID,
map/group, Stokes, spatial frame, WCS, pixel indexing, and support must be
explicit.  The accepted active lane is Stokes I; no polarimetric identity is
inferred.  Non-finite and missing states use declared validity metadata, not
undocumented numerical sentinels.

\section{Use-specific ensemble adequacy}

Let $Y^\epsilon$ be an observation, coadd, or filtered randomization product,
let $L_u$ be the fixed projection used by a named later consumer $u$, and let
$\mathcal V_u$ be that consumer's admitted validity domain.  Let
$\mathcal P_u^\star$, $\mu_u^\star$, and $T_u^\star$ denote an independently
authorized target law, mean, and covariance for that exact conditional use.
The following labels are deliberately nested only when their definitions make
them so:

\begin{align}
\text{distribution adequate for }u
&\iff
\Law(L_uY^\epsilon\mid\ThetaV,\mathcal V_u)
  =\mathcal P_u^\star,
\label{eq:distribution-adequacy}\\
\text{covariance adequate for }u
&\iff
L_u\mu_Y^\epsilon=\mu_u^\star\ \text{when required},\quad
L_uK_Y^\epsilon L_u^{\mathsf T}=T_u^\star,
\label{eq:covariance-adequacy}\\
\text{variance-scale adequate for }u
&\iff
\phi_u(L_uK_Y^\epsilon L_u^{\mathsf T})
=\phi_u(T_u^\star)
\label{eq:scale-adequacy}
\end{align}
for a predeclared scalar or diagonal functional $\phi_u$.  In an empirical
study the equalities are replaced only by owner-approved, preregistered
criteria.  This core sets no tolerance.  Scale adequacy does not imply
covariance adequacy; covariance adequacy does not imply distribution adequacy.

Before an ensemble can be an admissible input to any named use, it must also
have exact conditional-state identity, compatible units and shape, fixed and
matched operators, valid support, complete realized provenance, and adequate
assignment cardinality/dependence for that use.  The following consequences
remain within NOI-001:
\begin{itemize}
\item A direct pixelwise second-moment diagnostic requires only the exact
finite-stack identity in Equation~\eqref{eq:finite-stack-second-moment}; it
does not become a valid variance, weight, or signal-to-noise product.
\item A possible covariance input for NOI-002 or a fixed filter must satisfy
Equation~\eqref{eq:covariance-adequacy} on the joint pixel/operator domain;
matching only pixel diagonals is insufficient for aperture, matched-filter,
or extended-mode uses.
\item A possible global scale-calibration input may require only
Equation~\eqref{eq:scale-adequacy}, but that cannot authorize a spatial weight
model or confidence interpretation.
\item Blank positions or source-free regions can falsify failure to reproduce
a declared conditional ensemble, but they do not define an aperture estimator
or acceptance tolerance in this package.
\item The deterministic imprint in
Equation~\eqref{eq:signal-second-moment} can falsify a source-free-noise claim.
Its absence cannot prove distribution or covariance adequacy.
\end{itemize}

Whether a downstream sample covariance uses $R$, $R-1$, an effective-size
correction, shrinkage, a global rescaling, or another estimator; how it becomes
a weight or signal-to-noise value; and what count is required in production
are all outside these definitions and reserved to SCI-NOI-002.

\section{Analytic limits}

Any later implementation trace and fixture plan must recover these exact
limits before empirical interpretation:
\begin{enumerate}
\item $x_o=0$ implies $z_o^{(r)}=0$ and all moments vanish.
\item A deterministic sign assignment has zero randomization covariance by
Equation~\eqref{eq:deterministic-sign-law}, even though its uncentered second
moment can be nonzero.
\item One global centered sign gives
$M=K=(A_ox_o)(A_ox_o)^{\mathsf T}$ and rank at most one.
\item Independent centered coherence-unit signs give
\begin{equation}
  M=K=G_o(x_o)G_o(x_o)^{\mathsf T}
  =\sum_{a=1}^{K_o}g_{oa}(x_o)g_{oa}(x_o)^{\mathsf T}.
  \label{eq:independent-unit-limit}
\end{equation}
\item Uniform exact half balance obeys
Equation~\eqref{eq:within-realization-balance} and annihilates any map mode
with equal unit columns.
\item With $s_o=0$, deterministic-signal terms vanish, but the covariance is
still the conditional randomization covariance of one realized $n_o$, not
automatically the physical-noise covariance.
\item Oracle source removal gives Equation~\eqref{eq:perfect-signal-removal};
complete data overfit gives Equation~\eqref{eq:complete-overfit}.
\item Independent centered observation sign streams remove the $o\ne o'$
conditional randomization terms in Equation~\eqref{eq:coadd-covariance}; a
shared sign retains them.
\item Fixed linear coaddition and filtering obey
Equations~\eqref{eq:coadd-covariance} and
\eqref{eq:filter-propagation}; data-derived realization weights or filters do
not.
\item Complement assignments cancel empirical odd moments when paired but
repeat the same even second-moment pattern.
\end{enumerate}

\section{Proportional future tiny fixtures}

The following are specifications only.  They require no repository helper,
schema, verifier, reduction, or costly campaign and were not executed during
this freeze.  Any later execution requires the applicable Phase 2 and scope
authority.

\begin{enumerate}
\item \textbf{Two coordinates, identity map.}  Use $A=I_2$ and
$x=(a,b)^{\mathsf T}$.  Enumerate all signs to verify
Equations~\eqref{eq:ensemble-mean}--\eqref{eq:ensemble-covariance} for
$Q=I_2$, $Q=\one\one^{\mathsf T}$, and the $K=2$ balanced law
$Q=\bigl(\begin{smallmatrix}1&-1\\-1&1\end{smallmatrix}\bigr)$.

\item \textbf{Compact deterministic source.}  Use one output pixel,
$A=(1\;1)$, $s=(a,a)^{\mathsf T}$, and $n=0$.  The independent-unit,
global-coherent, and exactly balanced second moments are respectively
$2a^2$, $4a^2$, and $0$, verifying
Equation~\eqref{eq:signal-second-moment} and the coherence dependence.

\item \textbf{Physical versus randomization covariance.}  Use two physical
noise coordinates with symbolic
$\Sigma^{\mathrm{phys}}=\sigma^2
\bigl(\begin{smallmatrix}1&\rho\\\rho&1\end{smallmatrix}\bigr)$.
Independent signs remove the off-diagonal physical mode; one coherent sign
preserves it.  This verifies Equations~\eqref{eq:physical-map-covariance}--
\eqref{eq:preserved-physical-part} without selecting numerical values.

\item \textbf{Two-observation coadd.}  Use one pixel per observation and
fixed scalar coadd operators $H_1,H_2$.  Compare independent signs with a
shared sign to expose the exact cross term in
Equation~\eqref{eq:coadd-covariance}.

\item \textbf{Duplicate and complement census.}  For $K=2$, use the four
assignments $(++),(--),(+-),(-+)$.  Verify $U_{\mathrm{raw}}=4$ but
$U_{\pm}=2$, and verify that each complement pair has identical
$zz^{\mathsf T}$, as required by
Equations~\eqref{eq:raw-unique-count}--\eqref{eq:complement-unique-count}.

\item \textbf{Fixed filter propagation.}  Use a symbolic $2\times2$ matrix
$F$ and the two-coordinate covariance from fixture 1 to verify
$K_u=FK_yF^{\mathsf T}$ and to show that matching only
$\diag(K_y)$ does not determine the filtered variance.

\item \textbf{Residual extremes.}  Use the same $s,n$ with
$\widehat s=s$ and $\widehat s=s+n$.  Verify respectively the oracle-noise and
zero-residual limits in Equations~\eqref{eq:perfect-signal-removal} and
\eqref{eq:complete-overfit}.
\end{enumerate}

No fixture defines a scientific tolerance.  A realization count of 64 may be
considered only as an optional, resource-admitted high-count validation tier
under later authorization.  It is not a requirement, production default,
beammap expectation, or cap.

\section{Phase boundary and unresolved implementation questions}

This independent core establishes only the equations and contract above.  A
later separately authorized Phase 2 must bind them to the exact application
SHA by tracing, no more broadly than necessary:
\begin{itemize}
\item the insertion point and fixed conditioning state in
Equation~\eqref{eq:fixed-state};
\item the actual coherence partition and sign law in
Equations~\eqref{eq:coherence-lift}--\eqref{eq:sample-sign-covariance};
\item RNG, seed, balance, duplicate/complement, eligibility, and count state;
\item exact observation, coadd, validity, and filter operators in
Equations~\eqref{eq:map-operator}, \eqref{eq:fixed-coadd}, and
\eqref{eq:fixed-filter};
\item sequential/OpenMP reachability and scheduling independence;
\item realization persistence, units, shape, indexing, validity,
cardinality, and provenance; and
\item only the narrow fruit-loop state/interface needed to characterize
Equation~\eqref{eq:residual-ensemble}, without auditing SCI-FRUIT-001.
\end{itemize}
No scientific conclusion about implementation conformity, downstream
estimator validity, production use, or repair follows before that trace and
its separately governed evidence are complete.

\end{document}
